\begin{titlepage}
\thispagestyle{empty}
\vspace*{1.2cm}
{\sffamily\color{muted}\large EHB415 \textbullet\ Dönem Projesi}\par
\vspace{0.4cm}
{\sffamily\bfseries\fontsize{30}{34}\selectfont\color{acc} 5G Ağlarına Sensing\\[2pt] Entegrasyonu (ISAC)}\par
\vspace{0.3cm}
{\sffamily\Large Temelden uygulamaya öğretici}\par
\vspace{0.7cm}
{\color{acc!35}\rule{\textwidth}{1.4pt}}\par
\vspace{0.5cm}

\begin{minipage}[t]{0.62\textwidth}
\sffamily\normalsize
Bu belge, ekibin tamamının \emph{ne yaptığımızı, neden yaptığımızı ve nasıl yapacağımızı} aynı seviyede anlaması için yazıldı.
Soyut bir proposal'dan, çalıştırılabilir ve ölçülmüş bir başlangıca inmeyi hedefliyor.

\vspace{0.4cm}
\textbf{Bu belgede ne var?}
\begin{itemize}[leftmargin=1.2em]
\item Sorun ne, biz neyi düzeltmeye çalışıyoruz? (Bölüm~\ref{sec:buyuk})
\item Temel kavramlar: kaynak ızgarası, scheduler, OFDM radar (Bölüm~\ref{sec:kavram})
\item 3GPP standartları: ne söylüyorlar, biz neyi alıyoruz? (Bölüm~\ref{sec:3gpp})
\item Üç araç: Sionna, ns-3/5G-LENA, OCUDU (Bölüm~\ref{sec:arac})
\item OCUDU kodu: nerede ne var, nereye ekleme yaparız? (Bölüm~\ref{sec:ocudu})
\item Pilot görevler ve 4 kişilik görev dağılımı (Bölüm~\ref{sec:pilot})
\item Makale okuma rehberi ve Google Scholar sorguları (Bölüm~\ref{sec:okuma})
\end{itemize}
\end{minipage}\hfill
\begin{minipage}[t]{0.33\textwidth}
\vspace{0pt}
\begin{tcolorbox}[colback=accl,colframe=acc,boxrule=0.5pt,arc=3pt,left=6pt,right=6pt,top=5pt,bottom=5pt]
\sffamily\small
\textbf{Bu belgedeki sayılar nereden?}\par\vspace{2pt}
Kutular ve şekiller üç türdendir:
\begin{itemize}[leftmargin=1em,itemsep=1pt]
\item \textcolor{purp}{\textbf{Ölçülen bulgu}}: bu projenin kurulu araçlarında \emph{gerçekten çalıştırılıp} ölçüldü.
\item \textcolor{good}{\textbf{Pilot görev}}: sizin yapacağınız iş.
\item \textcolor{warn}{\textbf{Dikkat}}: tuzaklar ve doğrulanmamış yerler.
\end{itemize}
\end{tcolorbox}
\end{minipage}

\vfill
{\sffamily\small\color{muted}
Sürüm: 2026-10-17 \textbullet\ Araç sürümleri: Sionna 1.2.2, ns-3.48 + 5G-LENA v5.1, OCUDU 26.10.0\\
Kurulum dokümanları ve betikler: \texttt{kurulum/} klasörü \textbullet\ Okuma listesi: \texttt{makaleler/OKUMA\_LISTESI.md}}
\end{titlepage}

\tableofcontents
\clearpage

\section*{Nasıl okunmalı?}
\addcontentsline{toc}{section}{Nasıl okunmalı?}

Herkes \textbf{Bölüm~\ref{sec:buyuk}}'i (büyük resim) okur. Sonra rolüne göre:

\begingroup\sffamily\small
\begin{xltabular}{\textwidth}{@{}lXl@{}}
\toprule
\textbf{Rol} & \textbf{Önce oku} & \textbf{Pilot görevi} \\
\midrule
Herkes & Bölüm~\ref{sec:buyuk}, \ref{sec:kavram} ve~\ref{sec:pilot} & Pilot-1, 2, 3 (hepsi) \\
Standart ve metrik & Bölüm~\ref{sec:3gpp} & Bölüm~\ref{sec:pilot} \\
Sionna & Bölüm~\ref{sec:kavram}.4, \ref{sec:arac}.1 & Pilot-3 \\
ns-3 & Bölüm~\ref{sec:arac}.2 & Pilot-2 \\
OCUDU & Bölüm~\ref{sec:arac}.3, \ref{sec:ocudu} & Pilot-1 \\
\bottomrule
\end{xltabular}
\endgroup

\begin{ipucu}[Okurken şunu hatırla]
Bu proje bir \emph{yazılım yazma} projesi değil, bir \emph{ölçüm ve karşılaştırma} projesi. Hedefimiz yeni bir 5G sistemi yapmak değil;
mevcut açık kaynak araçlarla şu soruya kanıtlı bir cevap vermek: \textbf{sensing eklemek haberleşmeyi ne kadar etkiler?}
Bu yüzden bu belgenin her bölümü o soruya bir şey ekler.
\end{ipucu}
